\section{Axis 4: Sufficiency and Limits}
\label{sec:sufficiency}

The fourth axis concerns what a reduct \emph{cannot} guarantee.  Granular
computing finds a subset $R \subseteq \C$ that preserves the dependency
degree $\DEP(R) = \DEP(\C)$---but $\DEP(\C)$ itself may be far from 1.
Understanding the ceiling of what feature selection can achieve requires
results from classical learning theory that the GrC literature rarely cites.

%-----------------------------------------------------------------------
\subsection{Bayes Error Lower Bound}
\label{sec:bayes}
%-----------------------------------------------------------------------

The \emph{Bayes error} $R^* = \inf_f \mathbb{P}[f(X) \neq Y]$ is the
minimum achievable classification error over all measurable functions.
No feature set $R$, however well chosen, can reduce the empirical error
below $R^*$.

Cover and Hart~\cite{CoverHart1967} proved that the 1-NN rule achieves
error $R \leq R^*(2 - |\mathcal{D}|R^*/(|\mathcal{D}|-1))$ as $n\to\infty$;
in the binary case, $R \leq 2R^*(1-R^*)$, so $R \leq 2R^*$.  This result
implies that if a dataset has high inherent class overlap (high $R^*$), a
$k$-NN classifier on the reduct cannot do much better than twice the Bayes
error.

For finite samples, Drakopoulos~\cite{Drakopoulos1995} derives bounds on
classification error of the NN rule as a function of sample size and the
H\"older continuity of the class-conditional densities.  Wheat et
al.~\cite{Wheat2025BER} report Monte Carlo experiments showing that
practical BER estimators require at least 1000 samples per class to achieve
a 95\% confidence interval narrower than 5\%, and that accuracy degrades
rapidly as $d$ grows---a further endorsement of the concentration barrier
discussed in Section~\ref{sec:concentration}.

\paragraph{Implication for the dependency degree.}
When $\DEP(\C) < 1$, the dataset has a non-empty boundary region: some
objects are inherently ambiguous.  A reduct preserving $\DEP$ preserves
this ambiguity.  At a fixed radius $\delta$, $\DEP(\C)$ gives the fraction of training objects
whose neighbourhood falls entirely within one class (at that $\delta$); the remaining
$1-\DEP(\C)$ (boundary region) cannot be deterministically resolved by any
reduct at that $\delta$.  Unlike $R^*$, which is a population quantity,
$\DEP(\C)$ is an in-sample, $\delta$-dependent statistic: as $\delta\to 0$,
every granule becomes a singleton and $\DEP \to 1$, so $\DEP$ alone does not
bound generalisation error.  Accuracy on boundary objects is further
constrained by $R^*$.  The pair $(1-\DEP(\C),\,R^*)$ brackets the difficulty
of reduct-based classification from two complementary angles (in-sample
ambiguity and population-level irreducibility), yet the rough set literature
rarely reports or estimates either.

%-----------------------------------------------------------------------
\subsection{Class Overlap and Dataset Hardness}
\label{sec:overlap}
%-----------------------------------------------------------------------

Santos et al.~\cite{Santos2022overlap} provide a unified taxonomy of class
overlap measures, arguing that overlap (not imbalance per se) is the primary
driver of classification difficulty.  Their taxonomy distinguishes:
(a) \emph{feature overlap}: class conditional distributions share support;
(b) \emph{borderline overlap}: classes are locally mixed near the decision
boundary; (c) \emph{rare case overlap}: a few minority instances fall deep
in the majority region.

Han et al.~\cite{Han2024MFII} propose MFII, a complexity measure that
quantifies the joint effect of skewed distribution, overlap size, confusion
degree, and sub-cluster structure.  They also define ``value of
resolution'' (VoR) and ``degree of consistency'' (DoC) to evaluate whether
a complexity measure is discriminating: a measure with low VoR cannot
distinguish between datasets that respond differently to the same classifier.

The boundary region $\mathrm{BND}_R(d) = \UPP_R(d) \setminus \LOW_R(d)$
is a natural rough-set measure of class overlap.  The connection between
$|\mathrm{BND}_R(d)|/|\U| = 1 - \DEP(R)$ and the Santos--Han complexity
taxonomy has not been formalised.  Such a connection would allow GrC
practitioners to anticipate, before running reduction, whether a reduct
is likely to help or whether the dataset is in a regime where all methods
are near-optimal (i.e., near $R^*$).

%-----------------------------------------------------------------------
\subsection{Indirect and Weak Supervision}
\label{sec:weak}
%-----------------------------------------------------------------------

A reduct requires labelled training data to evaluate $\DEP(B)$.  Much
practical data arrives with indirect, noisy, or aggregate labels.

\paragraph{Learning from label proportions (LLP).}
In LLP, the learner observes only the class proportion within groups of
objects, not individual labels.  PAC-style analyses of LLP (see the survey by
Zhang et al.~\cite{Zhang2022PWS}, Section~4) suggest that learning requires
$O(1/\varepsilon^2)$ groups of size $O(\log(1/\delta))$ for
$(\varepsilon,\delta)$-accuracy; the sample complexity grows with the number
of classes but remains polynomially manageable.
A GrC-specific open problem is whether the dependency degree $\DEP(B)$
computed on aggregate label proportions converges to the clean-label
dependency degree as group size grows.

\paragraph{Partial label learning (PLL).}
In PLL, each object is annotated with a \emph{candidate set} of labels,
exactly one of which is correct.  Qian et al.~\cite{Qian2023InfoFusion}
combine granular ball computing with NRS to disambiguate candidate labels
before computing feature significance; their method adapts the rough set
reduct pipeline to the PLL setting with competitive accuracy on eight
benchmark datasets.  Wang et al.~\cite{Wang2018TNNLS} address the case
where some fraction of labels is randomly corrupted (not merely ambiguous),
providing an importance-reweighting strategy and consistency guarantees.

\paragraph{Programmatic weak supervision.}
Zhang et al.~\cite{Zhang2022PWS} survey programmatic weak supervision
(PWS), where labelling functions replace manual annotation.  The label model
in PWS estimates the accuracy of each labelling function; the resulting
probabilistic labels could, in principle, be used as a soft version of $d$
for computing $\DEP$.  This connection has not been explored in the GrC
literature.

\paragraph{Conformal prediction.}
Conformal prediction~\cite{Angelopoulos2023CP} provides finite-sample,
distribution-free prediction sets containing the true label with a
user-specified probability $1-\alpha$, without model or distributional
assumptions.  The positive region $\POS_R(d)$ offers a natural analogue: objects
in $\POS_R(d)$ receive singleton predictions (certain class membership) while
boundary objects are assigned multi-label sets.  This correspondence suggests
that reduct-based classification is implicitly performing a form of conformal
prediction with the nonconformity score defined by the neighbourhood
dependency.  Formalising this connection---including whether the positive-region coverage
provides distribution-free guarantees analogous to conformal prediction's
marginal coverage bound---has not been explored in the GrC literature.

\paragraph{Assessment of the sufficiency axis.}
The boundary between ``insufficient supervision'' and ``low dependency
degree'' is conceptually tight but notationally separate in the two
literatures.  A rough set reduct computed on noisy labels has a lower
$\DEP$ than the same reduct computed on clean labels; the gap between these
two $\DEP$ values is a natural measure of labelling noise.  This
interpretation is implicit in VPGB~\cite{Peng2022VPGB} and Qian et
al.~\cite{Qian2023InfoFusion} but has not been stated as a theorem.

Table~\ref{tab:sufficiency_synth} synthesises the four frameworks of this
section, showing what each bounds, whether GrC uses it, and where the gap
lies.

\begin{table*}[pos=!htbp]
\caption{Synthesis of sufficiency frameworks: what each bounds, its current
use in GrC, and the open gap.}
\label{tab:sufficiency_synth}
\centering\small
\begin{tabular}{p{3.0cm}p{3.8cm}p{3.0cm}p{4.5cm}}
\toprule
Framework & What it bounds & GrC adoption & Open gap \\
\midrule
Bayes error ($R^*$)
  & Population-level irreducible error
  & Rarely estimated; $\DEP$ used as proxy
  & Mapping $R^*$ to optimal purity $\tau$ (OP1) \\
\addlinespace
Class overlap (Santos--Han)
  & Dataset-level feature/border/rare-case overlap
  & Not connected to $\mathrm{BND}_R(d)$
  & Formal link between overlap profile and boundary region \\
\addlinespace
Weak/partial supervision
  & Label quality $\to$ $\DEP$ degradation
  & Implicit in VPGB, PLL-GBNRS
  & Tight bound on $|\DEP^* - \DEP^\eta|$ as function of noise rate (OP5) \\
\addlinespace
Conformal prediction
  & Marginal coverage at confidence $1{-}\alpha$
  & Not used in GrC
  & Whether POS coverage provides distribution-free guarantees \\
\bottomrule
\end{tabular}
\end{table*}
